% Level 1, Homework 10 -- My Art Show Plan (Milestone M1)   (input by homework.tex)
\homeworkweek{10}{My Art Show Plan}{15--20 minutes. This week's homework is your project plan (Milestone M1).}
\namedateline

\begin{core}[Finish your plan sheet (about 10 minutes)]
Use the \textbf{plan sheet} in your Project Brief (page 4). Make sure it has:
\begin{checklist}
  \item a \textbf{sketch} of your picture, with the parts labelled;
  \item your \textbf{question(s)}: 1 or 2 questions about choices (a number, a colour, day or night);
  \item your \textbf{function}: its name and its input slots, like \texttt{fish(x, y, colour)};
  \item a few words for each ingredient: my loop repeats\ldots, my \texttt{if} decides\ldots, my random surprise is\ldots, my variable remembers\ldots
\end{checklist}
Bring the plan sheet next week. You will build from it!
\end{core}

\begin{challenge}[Write your function on paper]
Write the recipe for your function. Start with the line below, then add the steps (penup, goto, pendown, then the drawing). If you have a computer, try it in your project file.
\par\medskip
\texttt{def} \rule{3cm}{0.4pt}\texttt{(x, y,} \rule{2cm}{0.4pt}\texttt{):}
\answerlines{6}
\end{challenge}

\begin{creative}[Colours and a title]
Choose 3 colours from the colour-blind-safe palette for your picture and give your artwork a title.
\begin{center}
\renewcommand{\arraystretch}{1.4}
\begin{tabular}{@{}l@{\qquad}l@{}}
\swatch{kcorange}\ orange \texttt{\#E69F00} & \swatch{kcblue}\ blue \texttt{\#0072B2}\\
\swatch{kcsky}\ sky blue \texttt{\#56B4E9} & \swatch{kcvermillion}\ vermillion \texttt{\#D55E00}\\
\swatch{kcgreen}\ green \texttt{\#009E73} & \swatch{kcpurple}\ purple \texttt{\#CC79A7}\\
\swatch{kcyellow}\ yellow \texttt{\#F0E442} & \swatch{kcblack}\ black \texttt{\#000000}\\
\end{tabular}
\end{center}
My 3 colours: \rule{8cm}{0.4pt}\par\smallskip
The title of my artwork: \rule{8cm}{0.4pt}
\end{creative}

\begin{selfcheck}
  \item My sketch shows what my program will draw.
  \item My questions are about choices, not about who someone is.
  \item My function has a name that says what it draws.
  \item I know where each of the six ingredients will go.
\end{selfcheck}
